(1) \(1-\cos u=2\sin^2(u/2)\) 使静态能量成为
\[
E=\frac12\int[u_x-2\sin(u/2)]^2dx+2\int u_x\sin(u/2)dx\ge8.
\]
第二项仅由边界的拓扑数决定，等号满足 \(u_x=2\sin(u/2)\)。积分得 \(\boxed{u_K=4\arctan e^{x-X}}\)，导数 \(2\operatorname{sech}(x-X)\)，直接代回 \(u_{xx}=\sin u\) 成立。拓扑荷 \((u(+\infty)-u(-\infty))/(2\pi)=1\)，静能恰为8。

(2) 由Lorentz不变性，\(u=4\arctan\exp[\gamma(x-vt-X)]\)、\(\gamma=(1-v^2)^{-1/2}\) 是精确移动孤子。此时 \(u_t=-vu_x\)，且 \(u_x=2\gamma\operatorname{sech}[\gamma(x-vt-X)]\)。代入能量和动量定义，令 \(s=\gamma(x-vt-X)\) 后积分 \(\operatorname{sech}^2s\)，得
\[
E=8\gamma,\qquad P=-\int u_tu_x\,dx=8\gamma v,\qquad E^2-P^2=64.
\]
因此孤子静质量8而不是随速度变化的数值；动量密度的负号与传播方向一致。

(3) 设 \(u=u_K+\epsilon\eta(x)e^{i\omega t}\)，保留一阶量，得 \([-\partial_x^2+\cos u_K]\eta=\omega^2\eta\)，其中 \(\cos u_K=1-2\operatorname{sech}^2(x-X)\)。令 \(A=\partial_x+\tanh(x-X)\)，有 \(A^\dagger A=-\partial_x^2+1-2\operatorname{sech}^2(x-X)\ge0\)。零模 \(\eta_0=\operatorname{sech}(x-X)/\sqrt2\) 归一化为1，伙伴算符 \(AA^\dagger=-\partial_x^2+1\)，故连续谱从 \(\omega^2=1\) 起，负特征值不可能出现。

(4) 将自由波 \(e^{ik(x-X)}\) 施以 \(A^\dagger\)，得到 \(\eta_k=[\tanh(x-X)-ik]e^{ik(x-X)}\)，且 \(\omega^2=k^2+1\)。两端只含正向行波，故反射振幅恒为0；透射幅 \(t(k)=(1-ik)/(-1-ik)\)，模为1，物理含义是无反射而有相移。零模与 \(\partial_Xu_K=-u_{K,x}\) 成比例，它表示平移对称性，不是虚频不稳定性。
